\section{Cuts, normal functions, and the reach of induction}
\label{sec:scope}
In this section $\mathcal O$ again means the original system below $\varepsilon_0$.
\subsection{Von Neumann identities as formulas about points}
Represent a possibly infinite collection by a given membership formula. Write $I_a(x)$ for $x<a$ with the usual guards. Then
\[
 I_a\subseteq I_b\ \longleftrightarrow\ a\le b,\quad
 I_a\cap I_b=I_{\min(a,b)},\quad I_{S(a)}=I_a\cup\{a\}.
\]
Here inclusion means pointwise implication. Transitivity proves inclusion when $a\le b$. If $b<a$, the point $b$ refutes it. The other formulas follow from comparison and immediacy of the successor. Thus even infinite represented ordinals admit a precise membership and inclusion reading in the arithmetic. For $a>0$, the pointwise union of its predecessors' cuts is $I_a$ if $a$ is a limit, and $I_c$ if $a=S(c)$. In the first case, $x<a$ implies $x<S(x)<a$, which gives an explicit predecessor containing $x$. In the second case $I_c$ itself occurs in the union and contains all its other summands. The empty case has empty union.

For a fixed PR sequence of codes $b(n)$, its union cut has the first-order membership formula $\exists n\ (x<b(n))$. If all $b(n)\le L$ and each $x<L$ has a witness $n$ with $x<b(n)$, the upper bound and the witness give the two inclusions proving that the union equals $I_L$. In general a formula describing a union need not yield a computable canonical name for it. For example let $b_e(n)$ be $1$ if a particular machine $e$ has halted by step $n$, and $0$ otherwise. These increasing-in-the-weak-sense sequences are uniformly PR. A total algorithm computing a canonical supremum for all of them would decide whether $e$ halts by testing the result against zero. If it halts the supremum is $1$, and if it never halts the supremum is zero. This is an external computability obstruction, already for a two-valued sequence.

\subsection{Two consequences of cofinality}
A fixed PR function $f$ on notations is called normal here if it is strictly increasing and has the following limit property. Whenever $z<f(L)$ at a positive limit, some specified or existentially given $y<L$ satisfies $z<f(y)$. For fixed $f$ its iterates $t_0=s$, $t_{n+1}=f(t_n)$ are PR. Suppose a positive limit $L$ is supplied with $t_n<L$ for all $n$ and with $\forall x<L\ \exists n\ x<t_n$. Then $L$ is the least fixed point of $f$ at or above $s$. Indeed $f(t_n)<f(L)$, so the tail sequence lies below $f(L)$. To see that the tail is cofinal, use $t_0<L$ and cofinality to get a term strictly above $t_0$, whose index is positive, and replace any witness at $t_0$ by that term. Thus $L\le f(L)$. If $L<f(L)$, continuity gives $y<L$ with $L<f(y)$. Choose $n$ with $y<t_n$. Then $f(y)<t_{n+1}<L$, a contradiction. Decidable comparison gives equality. For any fixed point $p\ge s$, numerical induction on the comparison $t_n\le p$ keeps all iterates below or equal to $p$, and so $L\le p$. The fixed-point argument uses the given limit name and cofinality. The towers supply both in the one-symbol extension with endpoint $E$.

Closed unbounded subsets of the whole original term order need not have nonempty intersection. In the tower sequence $T_0=1$, $T_{n+1}=[T_n]$, take separately the even and the odd indices. Both ranges are cofinal by the height bound and strict growth. Both membership tests are bounded. For $x$ of tree height $h$ only indices $n\le h+1$ could have $T_n\le x$. Each range has no limit point \emph{inside} $\mathcal O$. For $x>0$ choose the greatest listed point of that parity below $x$, when present. The interval $(p,S(x))$ then excludes all listed points except possibly $x$. If none is present, the open ray below $S(x)$ serves. Zero is isolated. The ranges are disjoint by strictness. This illustrates the uncountable regularity hypothesis in the article's club-intersection theorem.

There is also a useful cover statement at the first limit. Suppose a cover of $[0,\omega]$ has supplied choices of a covering open set at every point, by a PR selector (including an index for $\omega$), and the neighborhood data for the set chosen at $\omega$ give a tail bound $N$. Keep that one set and the choices at $0,\ldots,N-1$. Comparison of a point with $N$ proves that the list covers, using the tail for the remaining finite points and for $\omega$. This proves the finite-subcover form with the indicated data. Together with the neighborhood description it gives the witnessed one-point compactification of the discrete naturals.

\subsection{Hereditary bases and a sharp descent boundary}
For $b\ge2$ call a valid term $b$-admissible if every equal-exponent run, at every node, has length $<b$. Evaluate a word at base $b$ by
\[
 V_b([e_1,\ldots,e_k])=\sum_{j=1}^k b^{V_b(e_j)}.
\]
Define $H_b(n)$ by taking the base-$b$ digits of $n$ and using $H_b(j)$ as the exponent at each nonzero digit position $j$, with that digit as multiplicity. The expansion of zero is empty. When $n>0$, every used position $j$ has $j<n$ because $b^j\ge j+1$. Thus $H_b$ is defined by a course-of-values table on $n$, uniformly in $b$. Digits are found by bounded division and remainder, and the evaluation is a table on subtrees with $b$ a parameter.

\begin{theorem}[Base change and strict local descent]
For each $b\ge2$, $H_b$ and $V_b$ are inverse, order-preserving bijections between the naturals and the $b$-admissible terms. If $c\ge b$ and $n>0$, put $m=V_c(H_b(n))$. Then $m>0$ and
\[
       H_c(m)=H_b(n),\qquad H_c(m-1)<H_b(n).
\]
\end{theorem}
\begin{proof}
For admissible words, evaluation reflects the first CNF coefficient difference. To see this without assuming admissibility of any output, first prove this comparison lemma for pairs of admissible input trees. At each node the corresponding comparison assertion for child exponents is available by smaller-subtree induction. Apply it also to pairs of children from the same input, so distinct exponent values become distinct integer positions in the same order. If $j$ is the greatest position with a differing coefficient, the excess at $j$ is at least $b^j$ and the entire smaller-position part on either side is at most $\sum_{i<j}(b-1)b^i=b^j-1$. The last identity is ordinary induction. The empty-difference case gives equal values. This proves comparison including equality. For the induction one may fix the finite list of all subterms of the inputs and prove the comparison test for all pairs in it with maximum height at most $s$.

Now prove simultaneously by numerical table induction up to $n$ that $H_b(i)$ is valid and admissible and $V_b(H_b(i))=i$ for all $i\le n$. Distinct earlier digit positions have correct integer values, so their names are correctly ordered by the comparison lemma. Digits supply the permitted multiplicities, and ordinary base expansion proves the evaluation identity. For an admissible $u$, the trees $H_b(V_b(u))$ and $u$ have equal evaluations and hence are identical by equality reflection. These calculations give both inverse and order assertions. Finally $b$-admissibility implies $c$-admissibility and a nonempty word evaluates positively. Apply the inverse identity at $c$ and strict comparison to $m-1<m$.
\end{proof}
A special Goodstein step consists precisely of changing $b$ to $b+1$ and subtracting one from the positive value. The theorem proves decrease for each step and, by finite iteration, checks every supplied finite trace up to its first zero. Universal termination of all such sequences is unprovable even in PA, which interprets iPRA with the same PR definitions. This is the classical independence result of Kirby and Paris. See Rathjen's proof in the sources. Accordingly the stepwise decrease does not furnish a uniform termination proof in iPRA. The bound for every fixed $\omega^d$ in Section~\ref{sec:control} is a positive contrast with an explicit failure index.

\subsection{Why arbitrary presentations are different}
Here is an obstruction adapted from ProveIt's two-type promise note. Insert the number $0$ into the usual order of the positive integers by placing $x>0$ below it precisely when machine $e$ has not halted by time $x$. Positive-positive comparisons are unchanged. The positives placed below zero form an initial segment by monotonicity of bounded simulation, so we have a uniformly PR strict linear order. If $e$ never halts its order type is $\omega+1$. If $e$ first halts at $t$, zero comes immediately before $\max(1,t)$ in the positive list. The order type is $\omega$. A computable normalizer for presentations even with this promise, returning names with decidable equality to $\omega$, would decide the halting problem. This external argument explains why the algorithms take canonical notation, or explicit finite presentations with certificates. The article's proper class of all ordinals, unrestricted transfinite induction or recursion, uncountable cardinals, and stationary sets on regular uncountables have no literal quantifiers in the first-order numerical language. Finite-list escape ($S(\max A)$ above every listed term), witnessed suprema, the fixed-dimensional descent theorem and the two club families above are precise arithmetical counterparts.
